% Lec 03：量词、等号、Peano 算术与完备性
\documentclass[UTF8,a4paper,11pt]{ctexart}
\usepackage[a4paper,margin=2.3cm]{geometry}
\usepackage{amsmath,amssymb,mathtools}
\usepackage{enumitem}
\usepackage[hidelinks]{hyperref}
\usepackage{fancyhdr}
\setlist[itemize]{leftmargin=2.2em,itemsep=0.25em,topsep=0.35em}
\setlength{\parindent}{2em}
\setlength{\parskip}{0.35em}
\pagestyle{fancy}\fancyhf{}
\fancyhead[L]{离散数学与结构（I）}\fancyhead[R]{Lec 03}
\fancyfoot[C]{\thepage}\setlength{\headheight}{14pt}
\title{量词、等号、Peano 算术与完备性}
\author{zhangjunshuo}\date{2026 年 9 月}
\begin{document}
\maketitle
\begin{quote}
本讲从量词规则和等号推演进入 PA，再用 Henkin 构造说明一阶完备性，并由紧致性与 Löwenheim--Skolem 定理讨论非标准模型。
\end{quote}
\tableofcontents\clearpage
\section{量词规则与等号}


\subsection{一个完整的量词例子}


目标是

\[
\neg\forall x\neg\forall y\varphi(x,y)
\vdash
\forall y\neg\forall x\neg\varphi(x,y).
\]

证明的关键是两次使用全称引入时都避开开放假设中的特征变量：

\begin{verbatim}
1 | ¬∀x¬∀y φ(x,y)       PR
2 |   ∀x¬φ(x,a)          AS
3 |     ∀yφ(b,y)         AS
4 |     ¬φ(b,a)          ∀E 2
5 |     φ(b,a)           ∀E 3
6 |     ⊥                →E 4,5
7 |   ¬∀yφ(b,y)          →I 3-6
8 |   ∀x¬∀yφ(x,y)        ∀I 7   (b 不在开放假设中)
9 |   ⊥                  →E 8,1
10| ¬∀x¬φ(x,a)           →I 2-9
11| ∀y¬∀x¬φ(x,y)         ∀I 10  (a 不在开放假设中)
\end{verbatim}


第 10 行是由假设 $[\forall x\neg\varphi(x,a)]$ 导出矛盾后用 $\to I$ 得到否定；不能把它误写成 RAA。

\subsection{等号规则}


等号推演可用下列规则概括：

\begin{itemize}
\item 自反：$t=t$；
\item 对称与传递：由 $t_1=t_2$ 得 $t_2=t_1$，由 $t_1=t_2,t_2=t_3$ 得 $t_1=t_3$；
\item 函数相容：逐个参数相等时，函数值相等；
\item 公式代换：在原子式或更复杂公式中，把等项替换为等项。
\end{itemize}


最后一条是语义中的 Leibniz 原则：等号两边在所有项和谓词中具有相同的解释。
\clearpage
\section{Peano 算术}


\subsection{结构与理论要分开}


标准自然数结构写作

\[
\mathcal N=(\mathbb N;0,S,+,\cdot).
\]

说 $\mathcal N\models\mathrm{PA}$ 是说这个结构满足 PA 的所有公理；说 PA 是理论，则是说由一阶公理和推理规则得到的句子集合。结构不是理论，$\mathrm{PA}$ 也不等于 $\operatorname{Th}(\mathcal N)$。

\subsection{一组常用公理}


在语言 $(0,S,+,\cdot)$ 中，PA 包含：

\[
\begin{aligned}
&\forall x\,S(x)\ne0,\\
&\forall x\forall y\,(S(x)=S(y)\to x=y),\\
&\forall x\,(x+0=x),\qquad
\forall x\forall y\,(x+S(y)=S(x+y)),\\
&\forall x\,(x\cdot0=0),\qquad
\forall x\forall y\,(x\cdot S(y)=x\cdot y+x).
\end{aligned}
\]

最关键的是\textbf{一阶归纳公理模式}：对每个公式 $\varphi(x,\vec z)$，都有

\[
\bigl(\varphi(0,\vec z)\land\forall x(\varphi(x,\vec z)\to\varphi(S(x),\vec z))\bigr)
\to\forall x\varphi(x,\vec z).
\]

这里每个公式产生一个公理，不是一个能量化“所有公式”的单一一阶句子。

\subsection{做题时的区分}


\begin{itemize}
\item 集合论中的自然数归纳说某个集合是归纳集；
\item PA 的归纳模式说某个一阶公式在理论中可归纳；
\item 标准模型中的数学归纳法是语义事实。
\end{itemize}


三者结论相近，但层次不同，不能把其中一个直接当作另外两个的公理。
\clearpage
\section{Henkin 完备性}


\subsection{见证扩张}


设 $T$ 是一致的一阶理论。对每个存在句子 $\exists x\varphi(x)$ 加入新常元 $c_\varphi$ 和 Henkin 公理

\[
\exists x\varphi(x)\to\varphi(c_\varphi).
\]

逐层加入见证常元，确保前一层出现的每个存在句子都有见证；每层取演绎闭包。扩张保持一致，因为若新理论不一致，只会使用有限多个新常元和有限多个 Henkin 公理，可在原理论中用存在消去还原矛盾。

再用 Lindenbaum 引理把一致理论扩成极大一致集 $T^*$。经典极大一致性给出：对每个句子 $\theta$，恰有一个 of $\theta,\neg\theta$ 属于 $T^*$，并且联结词在 $T^*$ 中满足相应的真值条件。

\subsection{闭项模型}


论域取扩充语言的\textbf{所有闭项}按理论等价的商集：

\[
\Omega=\{t:t\text{ 是闭项}\}/\sim,
\qquad
t\sim u\Longleftrightarrow T^*\vdash t=u.
\]

定义 $[f(t_1,\ldots,t_n)]$ 为函数解释，定义关系

\[
R([t_1],\ldots,[t_n])\quad\Longleftrightarrow\quad R(t_1,\ldots,t_n)\in T^*.
\]

等号替换规则保证这些定义与代表元无关，因而是良定义的。

\subsection{真值引理与完备性}


按公式复杂度归纳证明真值引理：对任意把自由变量替换为闭项的实例，

\[
\mathcal M\models\varphi(t_1,\ldots,t_n)
\quad\Longleftrightarrow\quad
\varphi(t_1,\ldots,t_n)\in T^*.
\]

存在量词方向使用 Henkin 常元：若 $\exists x\varphi(x)\in T^*$，则 $\varphi(c_\varphi)\in T^*$，模型中有一个具体闭项见证。于是每个一致理论有模型；将它应用于 $\Gamma\cup\{\neg\varphi\}$，得到

\[
\Gamma\models\varphi\quad\Longrightarrow\quad\Gamma\vdash\varphi.
\]

这就是一阶逻辑的完备性方向；可靠性是另一方向 $\Gamma\vdash\varphi\Rightarrow\Gamma\models\varphi$。
\clearpage
\section{Löwenheim–Skolem 与非标准模型}


\subsection{紧致性制造大模型}


在 PA 的语言中加入一组常元 $\{c_s:s\in S\}$，并加入

\[
c_s\ne c_t\qquad(s\ne t).
\]

任意有限子集只要求有限多个常元两两不同，可以在标准自然数中解释；由紧致性，整个扩充理论有模型。因此 PA 有至少 $|S|$ 个元素的模型。取不可数的 $S$，就得到不可数的 PA 模型。

这说明一阶 PA 不能在同构意义下唯一刻画标准自然数；非标准模型中会出现标准部分以外的元素。

\subsection{向下 Löwenheim–Skolem}


若一个一阶理论使用可数语言，并且有无限模型，则向下 Löwenheim–Skolem 定理保证它有可数模型。对实闭域或集合论的相关一阶理论，也可得到可数模型；这里说的是理论在某个可数结构中成立，不是说外部的实数集合本身变成可数。

\subsection{Skolem 悖论的层次区别}


一个可数模型可以满足“没有从模型内部枚举出全部对象的函数”这类内部不可数性句子。模型内部的“可数”是指存在模型中的某个函数或集合；外部观察者看到的模型底层集合是否可数，是另一种元理论性质。两者不矛盾。

因此“可数模型满足不可数性”不是逻辑错误，而是内部量词与外部集合大小混用造成的表面悖论。

\end{document}
